% Supplementary submission materials for Current Science.
% Include these in the cover letter / submission portal.

\section*{Significance Statement (100 words)}
\noindent
Bearing fault diagnosis models trained on one operating condition often fail under load shifts. We show that the higher-order geometry of the reconstructed phase-space attractor, rather than the raw vibration waveform, is what transfers across a load shift. We introduce GA-DCHG, which screens which geometry descriptor is invariant per protocol, uses it to anchor hyperedge generation and cross-domain alignment, and validates the decision by test-time intervention. On two CWRU load-shift protocols GA-DCHG beats a 1D-CNN by 0.11--0.19 accuracy points. The anchor-screening result itself is protocol-dependent, a finding that should transfer to other phase-space-based diagnostic models.

\section*{Suggested Reviewers}
\noindent
1. Prof. Yongzhi Qu, School of Mechanical Engineering, Xi'an Jiaotong University, Xi'an, China. Email: yongzhiqu@xjtu.edu.cn. Research: graph neural networks for mechanical fault diagnosis.

2. Dr. Mohamed El Badaoui, Institut Pascal, Universite Clermont Auvergne, Clermont-Ferrand, France. Email: mohamed.elbadaoui@uca.fr. Research: phase-space methods and vibration signal processing for fault diagnosis.

3. Dr. Hui Ma, School of Engineering, RMIT University, Melbourne, Australia. Email: hui.ma@rmit.edu.au. Research: domain adaptation and transfer learning for rotating machinery fault diagnosis.

4. Prof. Te Han, School of Mechanical Engineering, Dalian University of Technology, Dalian, China. Email: hanspring@dlut.edu.cn. Research: hypergraph and graph learning for intelligent fault diagnosis.
